We include most relevant works to \fininteract and defer further related work in Appendix~\ref{app:related}.
%% ============================================================

\textbf{Financial QA benchmarks} such as FinanceBench~\citep{islam2023financebench}, FinQA~\citep{chen2021finqa}, and DocFinQA~\citep{reddy2024docfinqa} evaluate evidence-grounded answering and financial reasoning. Broader suites such as FinBen~\citep{xie2024finben} also cover generation, trading, and agent tasks. \fininteract focuses specifically on recovering a hidden intended reading through clarification, using paired default/intended interpretations. Appendix~\ref{app:related} compares the settings.

\textbf{Ambiguous QA and clarifying questions.} AmbigQA~\citep{min2020ambigqa} identifies plausible answers and rewrites questions to disambiguate them. CLAMBER~\citep{zhang2024clamber} evaluates ambiguity identification and clarification, while QuestBench~\citep{li2025questbench} studies information acquisition in reasoning tasks. \interactcomp~\citep{interactcomp2026} is our direct methodological ancestor, which we extend to finance with a domain-specific taxonomy, per-category targeting (AC@1), and verifiable financial evidence pairs.

\textbf{Multi-turn interaction and selective answering.} Large models often \emph{lose track} over multi-turn conversations~\citep{laban2026llms}, and calibrated systems frequently do better by answering \emph{selectively} or abstaining under ambiguity~\citep{cole2023selectively,kuhn2022clam}. Our findings echo this in the financial setting, where forcing extra clarification degrades accuracy, but \fininteract turns the observation into a controlled, per-category diagnosis of \emph{which} clarification an agent should issue and \emph{whether} issuing it resolves the query, instead of only measuring whether more turns help.
